\section{梯度下降真正难在步长}
\label{sec:09-Convex-Optimization/06-Optimization-Algorithms/01}

你有一个可微函数 \(f(x)\)，想找它的最小值。梯度指向函数增长最快的方向，所以负梯度是局部下降最自然的选择。迭代格式一行就能写完：

\[
x_{k+1}=x_k-\alpha_k\nabla f(x_k).
\]

看起来简单到不像是需要一整节讨论的问题。随便挑一个 \(\alpha_k\)，跑上几百步，不就完事了？

试一下就知道了。找一个最简单的二次函数 \(f(x)=\frac12 x^2\)，从 \(x_0=10\) 出发。取 \(\alpha=1\)，一步到零——完美。取 \(\alpha=0.01\)，每一步只挪当前位置的百分之一，100 步后离零点还远，收敛慢得像没动。取 \(\alpha=2\)，越过零点跳到 \(-10\)，下次跳回 \(10\)，来回震荡，目标值甚至比第一步前还大。

三个步长，三种命运。公式没变，方向没变，变天的是步长。你花一分钟写下迭代格式，可能花一整天调试步长。

这就是梯度下降的核心困境：方向由梯度免费提供，步长却决定了算法是收敛、爬行还是爆炸。更糟糕的是，同一个函数的不同区域，安全步长可能差好几个数量级——靠近最优点时应该更大还是更小？曲率大的方向应该谨慎还是大胆？直觉在这些问题上经常给出相反的答案。

\subsection{光滑性把``足够小''变成可计算条件}

问题在于没法凭空判断一个步长是否安全。太小浪费计算；太大会越过极小点，甚至让目标值反而增大。需要的是一个能事先计算的上界——不靠试错，而靠函数本身的某种性质。

关键线索是函数的``光滑程度''。梯度变得太快的地方，线性近似失效就快，能用的步长就小。如果梯度在单位距离内的变化不超过某个值，你就有了一把尺子。把这种直觉正式化：假设梯度是 \(L\)-Lipschitz 的，

\[
\norm{\nabla f(x)-\nabla f(y)}
\le L\norm{ x-y}.
\]

这个条件说的是：梯度不会振荡得太剧烈——两点之间梯度差异有上界，且与距离成比例。\(L\) 是曲率的上限，控制了一阶 Taylor 近似的误差。如果你知道了 \(L\)，就同时知道了线性近似的可靠半径。

从微积分基本定理出发，沿路径积分：

\[
f(y)-f(x)=\int_0^1\nabla f(x+t(y-x))^{\trans}(y-x)\,dt.
\]

被积函数里的梯度 \(\nabla f(x+t(y-x))\) 不是 \(\nabla f(x)\) 本身，而是沿路径各点的梯度。把 \(\nabla f(x)\) 拆出来：

\[
\begin{aligned}
f(y)-f(x) &= \nabla f(x)^{\trans}(y-x) \\
&\quad + \int_0^1\bigl[\nabla f(x+t(y-x))-\nabla f(x)\bigr]^{\trans}(y-x)\,dt.
\end{aligned}
\]

第一项是局部线性近似——这就是梯度下降所依赖的全部信息。第二项是近似误差，它的被积函数由 Lipschitz 条件控制：路径上任意一点到 \(x\) 的梯度变化不超过 \(L\) 乘以该点到 \(x\) 的距离。做积分计算，得到

\[
f(y)
\le
f(x)+\nabla f(x)^{\trans}(y-x)
+\frac L2\norm{ y-x}^2.
\]

这个不等式是一个全局的上方控制（global upper bound）：在任意一点 \(y\)，函数值不会超过以 \(x\) 为中心的某个二次函数的值。二次项系数 \(L/2\) 恰是最大曲率的一半。如果实际曲率比 \(L\) 小得多，这个界就偏松——但在缺乏更精确信息时，它是最可靠的保守估计。

现在代入具体的梯度下降步 \(y=x-\alpha\nabla f(x)\)。经过代入和整理：

\[
f(x-\alpha\nabla f(x))
\le
f(x)
-\alpha\left(1-\frac{L\alpha}{2}\right)
\norm{\nabla f(x)}^2.
\]

看括号因子 \((1-L\alpha/2)\)。它若为正，整个惩罚项非负，右侧严格小于 \(f(x)\)——函数值确定下降。让它为正的条件就是

\[
0<\alpha<\frac2L.
\]

取 \(\alpha=1/L\) 时，括号等于 \(1/2\)，每一步至少减少 \(\frac{1}{2L}\norm{\nabla f(x)}^2\)——下降量被梯度和光滑常数联合决定。

注意推导的走向：\(L\) 是梯度变化的最大速率，即曲率上界。步长上限不由函数值大小决定，而由曲率决定。梯度很大的区域完全可能允许大步长——只要那里曲率小。真正危险的是在曲率大的方向上迈出过长的一步，使二次上界失准。这也是为什么在高维中，各向异性（不同方向曲率差异大）比单纯的梯度绝对值更让人头疼。

\subsection{狭长二次碗为什么让梯度下降曲折}

上一段给出了安全步长上界，但安全不等于高效。考虑一个具体的二维函数，它能清楚地暴露问题：

\[
f(x,y)=\frac12x^2+50y^2.
\]

两个方向的曲率差距悬殊：\(x\) 方向曲率为 \(1\)，\(y\) 方向曲率为 \(100\)。Hessian 特征值为 \(1\) 和 \(100\)，光滑常数为 \(L=100\)。

现在做梯度下降。步长必须由最陡峭的 \(y\) 方向决定——\(\alpha=1/L=0.01\)——否则 \(y\) 坐标会越过谷底振荡发散。但这个步长在平缓的 \(x\) 方向上，每次只前进两轮精度：从 \(x_0=10\) 到零点，梯度在 \(x\) 方向上的分量只有 \(x\) 本身，每步只前进 \(0.01\cdot x\)，几百步后才能接近零点。

等高线是一组狭长椭圆，长短轴之比约为 \(10:1\)。梯度总是垂直于等高线——这意味着梯度方向近乎与长轴（\(x\) 方向）垂直。所以每一步走的方向，都不是指向远处的椭圆中心，而是近乎横跨谷地的。下一步从新位置的等高线再出发，又被推向对面谷壁。迭代轨迹在谷地两侧左右摆动，像之字形锯齿，缓慢地沿长轴朝中心蠕动。

\begin{center}
\begin{tikzpicture}[x=1cm,y=1cm,>=stealth]
  \draw[gray!35] (0,0) ellipse [x radius=4.2,y radius=1.35];
  \draw[gray!45] (0,0) ellipse [x radius=3.2,y radius=1.02];
  \draw[gray!55] (0,0) ellipse [x radius=2.2,y radius=0.70];
  \draw[gray!65] (0,0) ellipse [x radius=1.2,y radius=0.38];
  \fill (0,0) circle (2pt) node[below] {最小点};
  \draw[->,blue!70!black,thick]
    (4.0,1.05) -- (3.25,-0.88) -- (2.55,0.70) --
    (1.90,-0.52) -- (1.30,0.36) -- (0.82,-0.23) --
    (0.42,0.11) -- (0.08,-0.02);
  \node[above right] at (4.0,1.05) {\(x_0\)};
  \node[above] at (2.8,1.5) {陡峭方向来回穿越};
  \draw[->,gray!80] (2.7,1.35) -- (2.58,0.78);
  \node[below] at (2.7,-1.45) {平缓方向推进很慢};
\end{tikzpicture}

\emph{安全步长由最陡方向限制，于是轨迹横跨狭谷反复折返，只能缓慢地沿长轴接近中心。}
\end{center}

这就是大条件数的诅咒。对 \(m\)-强凸、\(L\)-光滑的函数，定义条件数

\[
\kappa=L/m.
\]

固定步长 \(1/L\) 时，收敛是几何的，但收缩因子大约为 \(1-m/L = 1-1/\kappa\)。\(\kappa=100\) 时每次迭代误差大约乘 \(0.99\)，一个数量级需要大约 230 步。\(\kappa=1000\) 时需要约 2300 步。推而广之，达到 \(\varepsilon\) 精度需要的步数约正比于 \(\kappa\log(1/\varepsilon)\)。

换个角度看同一件事。如果把坐标伸缩一下——令 \(z=\sqrt{100}\,y\)——\(y\) 方向被压缩，\(f\) 在新坐标中变成 \(f(x,z)=\frac12x^2+\frac12z^2\)，一个完美的圆。同样的固定步长在所有方向均匀推进，一步就能收敛至零点附近。改变的不是算法，是坐标系。这就是预条件（preconditioning）的核心想法：不换步长，换一个让问题不那么扁的坐标系。条件数降下来，梯度下降自然会走得快。

\subsection{\texorpdfstring{不知道 \(L\) 时怎样选择步长}{不知道 L 时怎样选择步长}}

前面的分析假设你已知道全局的 \(L\)。真实问题中精确的 \(L\) 经常无从得知——你可能只知道一些局部的或近似的估计。必须猜一个吗？

有更好的办法。回溯线搜索（backtracking line search）从一个候选步长出发，逐步缩小，直到函数出现了``足够''的下降。具体检查的是 Armijo 条件：

\[
f(x+\alpha p)
\le
f(x)+c\alpha\nabla f(x)^{\trans}p,
\qquad0<c<1.
\]

不等号右边是一个线性函数——由当前点的函数值和方向导数构成的一条直线。沿着方向 \(p\) 走 \(\alpha\) 距离后，如果实际函数值超出了这条直线（意味着线性近似已经严重不准），就说明步子太大。把 \(\alpha\) 乘以一个小于 1 的因子（典型值如 \(0.5\)），再试一次。这种收缩会持续到条件成立为止。

对负梯度方向 \(p=-\nabla f(x)\)，右侧是 \(f(x)-c\alpha\norm{\nabla f(x)}^2\)。条件明确要求目标下降一个与步长和梯度平方成正比的量。常数 \(c\) 通常取一个小值（如 \(10^{-4}\)），放宽要求，避免因浮点误差或轻微的非线性就让步长被无限收缩。

回溯不沿射线找精确极小值。精确线搜索本身对非二次函数的代价可能比原问题还高——你以为是在省步数，其实每步花的时间可能够做十次回溯。在有噪声的目标上，精确线搜索更是毫无意义：噪声会把``精确最小值''变成一个随机变量。

实现中还有几处经验的弯弯绕绕：

\begin{itemize}
\tightlist
\item
  初始步长沿用上一次成功的值。如果上一步附近的地形和当前位置相似，沿用步长比从头猜快得多。但如果地形陡变（接近边界、进入高曲率区），沿用的步长可能过大。
\item
  函数值评估可能含浮点噪声或 mini-batch 随机波动，需要在 Armijo 条件中留出余量，或采用非单调线搜索。
\item
  有约束时不只是沿直线走——可能还要投影回可行域，Armijo 条件也需要在约束框架中重新解释。
\item
  目标函数若昂贵（如基于全量数据的 loss），回溯次数需要硬性限制——有时宁可接受次优步长。
\item
  下降条件用的是全批目标还是随机估计？若是 mini-batch 估计，Armijo 条件的函数值带有方差，确定性理论需要重新建立。
\end{itemize}

\subsection{收敛速度依赖额外结构}

光滑性单独给出的是每一步下降的量。但``每步下降 \(\frac{1}{2L}\norm{\nabla f(x_k)}^2\)''本身不告诉你离终点还有多远——如果梯度一直很大，说明还在陡坡上；如果梯度已经很小，也可能差着十万八千里（非凸情形尤其如此）。

要把每步下降的累积翻译成``离最优解还有多远''，需要凸性连接局部和全局。凸性保证：函数值误差不超过内积 \(\nabla f(x_k)^{\trans}(x_k-x^\star)\)，而梯度范数和距离又通过光滑性和强凸性互相制约。

对 \(L\)-光滑的凸函数，固定步长 \(1/L\) 的梯度下降给出

\[
f(x_k)-f^\star=O(1/k).
\]

即每轮迭代将误差缩减一个大致与 \(1/k\) 成比例的量。十轮后误差约是初始的 \(1/10\)，百轮后约 \(1/100\)。不是特别快，但在无更强假设下无法再改善。

如果函数还是 \(m\)-强凸的——函数值下界不是平面，而是被一个二次函数从下方撑住——收敛大幅加速：

\[
f(x_k)-f^\star
\le
\left(1-\frac mL\right)^k
\bigl(f(x_0)-f^\star\bigr).
\]

每次迭代将误差乘以常数因子 \((1-m/L)\)。十轮后误差约是 \((1-1/\kappa)^{10}\)——慢还是快，完全看 \(\kappa\)。

术语上有个容易踩坑的地方：``线性收敛''指的是误差每一步乘以固定的小于 1 的因子，导致对数误差与迭代次数成线性关系——所以叫``线性''。它不是指误差随 \(k\) 线性下降（即 \(O(1/k)\)）。看着公式推导才不会弄混这两个完全不同的速率概念。

光滑凸问题的 \(O(1/k)\) 是否最优？不是。Nesterov 证明了一阶方法的下界是 \(O(1/k^2)\)，并构造了达到这一下界的加速方法。加速方法不改变每一步的计算成本——仍只需一次梯度求值——仅通过精心设计的带动量的辅助序列，把速率从 \(O(1/k)\) 推到 \(O(1/k^2)\)。它揭示了一阶方法的设计空间比表面看上去深得多：梯度能提供的不只是局部下降方向，还能像二阶方法那样编码全局几何信息——如果你知道怎么组合它。

但加速序列对参数敏感。在噪声环境中（如随机梯度、mini-batch），辅助序列的动量可能放大方差而非加速收敛。实践中经常需要重启策略：每隔固定步数重置动量，牺牲部分理论速率换鲁棒性。

\subsection{什么时候可以停止}

相邻两次迭代几乎不动，可能因为接近最优，也可能因为步长太小、困在了平坦区。函数值不再下降，可能是到了底线，也可能是因为步长被过分保守的线搜索压到了近乎无效的程度。``看起来收敛了''不等于``已经收敛了''。

更有意义的停止证据包括：

\begin{itemize}
\tightlist
\item
  无约束凸问题：梯度范数——一阶必要条件明确说最优点梯度为零，且凸性保证小梯度隐含函数值接近最优；
\item
  约束问题：投影梯度范数或 KKT 残差——这是在约束流形上的最优性一阶条件；
\item
  已知强凸常数时：从梯度范数或函数值反推的误差上界——但常数估计本身就难；
\item
  原始-对偶间隙：同时求解原问题和对偶问题，两者函数值的差距是一个可计算的全局次优度证书；
\item
  更实际的尺子：验证集表现不再改善、统计上误差已低于数据噪声的量级、或应用所需的决策精度已经满足。
\end{itemize}

优化误差不必远小于数据噪声和建模误差。把梯度降到机器精度可能浪费大量计算，还可能把病态方向上的数值误差放大到主宰整个更新。在机器学习中，通常验证集性能在优化误差还相当大时就已饱和——此时继续压低训练损失反而可能损害泛化。

最后，注意一条容易被忽略的线索：这一整节的下降引理推导中，凸性一次都没有被用到——只用了光滑性。凸性的职责在后面：它负责把局部下降翻译成全局函数值误差保证，把``每步都降了一点''垒成``离目标确实近了一步''。

\hyperref[ch:078]{``非凸与随机优化：能下降不等于找到全局解''}保留了同一条下降不等式，却把结论降级为驻点复杂度，并进一步引入鞍点、PL 条件与随机梯度噪声。失去凸性后，下降仍在，但下降能否通向足够好的解，是另一个独立的故事。
